\documentclass[UTF8, 12pt, fontset=fandol, oneside]{ctexart}
\usepackage{researchnotes}

% 保留分册独立编译时的章号前缀。
\renewcommand{\thesection}{8.\arabic{section}}

\title{第八章\quad 广义积分}
\author{}
\date{}

\begin{document}

\maketitle

在定积分中讨论的函数必须是定义在有限区间上的有界函数, 在本章中我们将定积分作两个方面的推广: 一是将有限区间推广到无穷区间, 二是将有界函数推广到无界函数. 对于前者的积分, 我们称之为无穷积分, 而将后者的积分称为瑕积分. 通常将无穷积分和瑕积分统称为广义积分。

\section{无穷积分的基本概念与性质}

在本节中我们研究无穷积分. 首先我们来考查一个火箭升空的例子. 假设人们在地球上某个地方向天空发射一个质量为 $m$ 的火箭, 其目的是探索太阳系中某个遥远的星球. 设地球半径为 $R$, 质量为 $M$, 则由万有引力定律知, 火箭从地球表面发射升至高度 $X$, 它克服地球引力所做的功为
\[
  W(X)=\int_R^{R+X}G\frac{Mm}{r^2}\mathrm dr,
\]
其中 $G=\dfrac{R^2g}{M}$, 而 $g$ 为大家熟悉的重力加速度, 即 $g\approx9.8\mathrm{m}/\mathrm{s}^2$. 对任意固定的 $X$, 我们有
\[
  W(X)=mgR^2\left(\frac1R-\frac1{R+X}\right).
\]
火箭逐步摆脱地球引力的过程可以看成是 $X\to+\infty$ 的变化过程. 因此, 从理论上来看, 火箭要摆脱地球引力所做的功应该为
\[
\begin{aligned}
  \lim_{X\to+\infty}W(X)
  &=\lim_{X\to+\infty}\int_R^{R+X}G\frac{Mm}{r^2}\mathrm dr\\
  &=\lim_{X\to+\infty}mgR^2\left(\frac1R-\frac1{R+X}\right)=mgR.
\end{aligned}
\]
利用上述分析, 我们可以求出所谓的第二宇宙速度, 即火箭具有的最小初速度, 使得它能够飞离地球的引力范围. 事实上, 设火箭向上发射时的初速度为 $v_0$, 它得到的动能为 $\dfrac12mv_0^2$. 当它超过 $\lim\limits_{h\to+\infty}W(h)=mgR$ 时, 则火箭可以脱离地球的引力. 因此由能量守恒定律, 有
\[
  \frac12mv_0^2=mgR,
\]
取 $g=9.81\mathrm{m}/\mathrm{s}^2,R=6371\ \mathrm{km}$, 代入得
\[
  v_0\approx11.2\ \mathrm{km}/\mathrm{s}.
\]

现在我们只关注火箭在升空的过程中克服地球引力做功所涉及的积分问题. 在火箭达到的每一个高度 $X$, 我们都面临一个有限区间的积分 $\int_R^X G\dfrac{Mm}{r^2}\mathrm dr$, 该积分对固定的 $X$ 是一个定积分, 但由于 $X$ 不断变大, 实际我们将面临的是考查当 $X\to+\infty$ 时, 积分 $\int_R^X G\dfrac{Mm}{r^2}\mathrm dr$ 的变化趋势. 我们把这个考查过程记为 $\int_R^{+\infty}G\dfrac{Mm}{r^2}\mathrm dr$, 并称之为无穷积分. 为此给出下面的定义。

\noindent\textbf{定义 8.1.1}\quad 设函数 $f(x)$ 在 $[a,+\infty)$ 上有定义, 并且对于 $\forall X\in(a,+\infty)$, 在 $[a,X]$ 上可积. 如果极限
\[
  \lim_{X\to+\infty}\int_a^X f(x)\mathrm dx
\]
存在, 则称无穷积分 $\int_a^{+\infty}f(x)\mathrm dx$ 收敛, 或称函数 $f(x)$ 在 $[a,+\infty)$ 上可积, 并记
\[
  \int_a^{+\infty}f(x)\mathrm dx=\lim_{X\to+\infty}\int_a^X f(x)\mathrm dx;
\]
如果极限 $\lim\limits_{X\to+\infty}\int_a^X f(x)\mathrm dx$ 不存在, 则称无穷积分 $\int_a^{+\infty}f(x)\mathrm dx$ 发散。

从上述定义可以看出: 记号 $\int_a^{+\infty}f(x)\mathrm dx$ 仅代表了我们考虑 $f(x)$ 在 $[a,+\infty)$ 上的积分这件事情, 只有当无穷积分 $\int_a^{+\infty}f(x)\mathrm dx$ 收敛时, 它才是一个有限实数。

类似地我们也可以讨论一个函数在 $(-\infty,b]$ 上的积分问题。

\noindent\textbf{定义 8.1.2}\quad 设函数 $f(x)$ 在 $(-\infty,b]$ 上有定义, 并且对于 $\forall X\in(-\infty,b)$, 在区间 $[X,b]$ 上可积. 如果极限
\[
  \lim_{X\to-\infty}\int_X^b f(x)\mathrm dx
\]
存在, 则称无穷积分 $\int_{-\infty}^b f(x)\mathrm dx$ 收敛, 或称函数 $f(x)$ 在 $(-\infty,b]$ 上可积, 并记
\[
  \int_{-\infty}^b f(x)\mathrm dx=\lim_{X\to-\infty}\int_X^b f(x)\mathrm dx;
\]
如果极限 $\lim\limits_{X\to-\infty}\int_X^b f(x)\mathrm dx$ 不存在, 则称无穷积分 $\int_{-\infty}^b f(x)\mathrm dx$ 发散。

\noindent\textbf{定义 8.1.3}\quad 设函数 $f(x)$ 在 $(-\infty,+\infty)$ 上有定义, 且在任何的闭区间 $[a,b]$ 上可积. 任取 $c\in\mathbb R$, 若无穷积分 $\int_{-\infty}^c f(x)\mathrm dx$ 与 $\int_c^{+\infty}f(x)\mathrm dx$ 都收敛, 则称无穷积分 $\int_{-\infty}^{+\infty}f(x)\mathrm dx$ 收敛, 并记
\[
  \int_{-\infty}^{+\infty}f(x)\mathrm dx
  =
  \int_{-\infty}^c f(x)\mathrm dx
  +
  \int_c^{+\infty}f(x)\mathrm dx;
\]
若无穷积分 $\int_{-\infty}^c f(x)\mathrm dx$ 和 $\int_c^{+\infty}f(x)\mathrm dx$ 中至少有一个发散, 则称无穷积分 $\int_{-\infty}^{+\infty}f(x)\mathrm dx$ 发散。

由于我们在定义 8.1.3 中假定了 $f(x)$ 在任何有限闭区间均可积, 读者易于验证, $\int_{-\infty}^{+\infty}f(x)\mathrm dx$ 收敛与否与 $c$ 的选取无关; 当 $\int_{-\infty}^{+\infty}f(x)\mathrm dx$ 收敛时, 其值也与 $c$ 的选取无关。

由无穷积分敛散性的定义我们可以看出: 若函数 $f(x)$ 在 $[a,+\infty)$ 上有原函数 $F(x)$, 并形式地记 $F(+\infty)=\lim\limits_{x\to+\infty}F(x)$, 则有
\[
  \int_a^{+\infty}f(x)\mathrm dx=F(x)\big|_a^{+\infty}=F(+\infty)-F(a).
\]
同样, 若 $f(x)$ 在 $(-\infty,b]$ 上有原函数 $G(x)$, 记 $G(-\infty)=\lim\limits_{x\to-\infty}G(x)$, 则
\[
  \int_{-\infty}^b f(x)\mathrm dx=G(x)\big|_{-\infty}^b=G(b)-G(-\infty).
\]
类似地, 若 $f(x)$ 在 $(-\infty,+\infty)$ 上有原函数 $H(x)$, 则
\[
  \int_{-\infty}^{+\infty}f(x)\mathrm dx=H(x)\big|_{-\infty}^{+\infty}=H(+\infty)-H(-\infty).
\]

上述的公式给出了判定无穷积分是否收敛以及计算收敛无穷积分的值的方法。

由定积分几何意义, 不难理解无穷积分 $\int_a^{+\infty}f(x)\mathrm dx$, $\int_{-\infty}^b f(x)\mathrm dx$ 和 $\int_{-\infty}^{+\infty}f(x)\mathrm dx$ 的几何意义. 例如, 当 $f(x)$ 是 $[a,+\infty)$ 上的连续函数时, $\int_a^{+\infty}f(x)\mathrm dx$ 的几何意义是图 8.1.1 中阴影部分面积的代数和。

\noindent\textbf{例 8.1.1}\quad 证明无穷积分 $\int_{-\infty}^{+\infty}\dfrac{\mathrm dx}{1+x^2}$ 收敛并求其值。

\noindent\textbf{解}\quad 由于 $\arctan x$ 是 $\dfrac1{1+x^2}$ 在 $(-\infty,+\infty)$ 上的一个原函数, 因此我们有
\[
  \lim_{X\to-\infty}\int_X^0\frac{\mathrm dx}{1+x^2}
  =
  0-\lim_{X\to-\infty}\arctan X
  =
  \frac{\pi}{2},
\]
\[
  \lim_{X\to+\infty}\int_0^X\frac{\mathrm dx}{1+x^2}
  =
  \lim_{X\to+\infty}\arctan X
  =
  \frac{\pi}{2}.
\]
由无穷积分的定义知 $\int_{-\infty}^{+\infty}\dfrac{\mathrm dx}{1+x^2}$ 收敛, 并且有
\[
  \int_{-\infty}^{+\infty}\frac{\mathrm dx}{1+x^2}
  =
  \arctan x\big|_{-\infty}^{+\infty}
  =
  \frac{\pi}{2}-\left(-\frac{\pi}{2}\right)=\pi.
\]

\noindent\textbf{例 8.1.2}\quad 讨论无穷积分 $\int_1^{+\infty}\dfrac{\mathrm dx}{x^p}$ 的敛散性, 其中 $p\in\mathbb R$。

\noindent\textbf{解}\quad 当 $p=1$ 时,
\[
\begin{aligned}
  \int_1^{+\infty}\frac{\mathrm dx}{x}
  &=
  \lim_{X\to+\infty}\int_1^X\frac{\mathrm dx}{x}
  =
  \lim_{X\to+\infty}\ln x\big|_1^X\\
  &=
  \lim_{X\to+\infty}\ln X=+\infty;
\end{aligned}
\]
当 $p\ne1$ 时,
\[
\begin{aligned}
  \int_1^{+\infty}\frac{\mathrm dx}{x^p}
  &=
  \lim_{X\to+\infty}\int_1^X\frac{\mathrm dx}{x^p}
  =
  \lim_{X\to+\infty}\frac{x^{-p+1}}{1-p}\bigg|_1^X\\
  &=
  \lim_{X\to+\infty}\frac{X^{-p+1}-1}{1-p}
  =
  \begin{cases}
    \dfrac1{p-1}, & p>1,\\
    +\infty, & p<1.
  \end{cases}
\end{aligned}
\]
因此无穷积分 $\int_1^{+\infty}\dfrac{\mathrm dx}{x^p}$ 当 $p>1$ 时收敛, 当 $p\leqslant1$ 时发散。

\noindent\textbf{例 8.1.3}\quad 讨论无穷积分 $\int_2^{+\infty}\dfrac{\mathrm dx}{x\ln^p x}$ 的敛散性, 其中 $p\in\mathbb R$。

\noindent\textbf{解}\quad 当 $p=1$ 时,
\[
\begin{aligned}
  \int_2^{+\infty}\frac{\mathrm dx}{x\ln x}
  &=
  \lim_{X\to+\infty}\int_2^X\frac{\mathrm dx}{x\ln x}\\
  &=
  \lim_{X\to+\infty}(\ln\ln X-\ln\ln2)=+\infty.
\end{aligned}
\]
当 $p\ne1$ 时,
\[
\begin{aligned}
  \int_2^{+\infty}\frac{\mathrm dx}{x\ln^p x}
  &=
  \lim_{X\to+\infty}\int_2^X\frac{\mathrm dx}{x\ln^p x}
  =
  \lim_{X\to+\infty}\frac{\ln^{-p+1}x}{1-p}\bigg|_2^X\\
  &=
  \lim_{X\to+\infty}\frac{\ln^{1-p}X-\ln^{1-p}2}{1-p}
  =
  \begin{cases}
    \dfrac1{(p-1)\ln^{p-1}2}, & p>1,\\
    +\infty, & p<1.
  \end{cases}
\end{aligned}
\]
因此无穷积分 $\int_2^{+\infty}\dfrac{\mathrm dx}{x\ln^p x}$ 当 $p>1$ 时收敛, 当 $p\leqslant1$ 时发散。

类似地, 我们可以证明: 无穷积分
\[
  \int_{e^2}^{+\infty}
  \frac{\mathrm dx}{x(\ln x)(\ln\ln x)^p}
\]
当 $p>1$ 时收敛, 当 $p\leqslant1$ 时发散。

从无穷积分的定义出发, 我们可以容易地证明下面的关于无穷积分的换元公式与分部积分公式。

设函数 $f(x)$ 在 $[a,+\infty)$ 上有定义, 且对于 $\forall X>a$, 在区间 $[a,X]$ 上可积, 再设函数 $\varphi(t)$ 在区间 $[\alpha,\beta)(\beta$ 可以是 $+\infty)$ 上连续可微, 严格单调上升, 并且满足
\[
  a=\varphi(\alpha)\leq \varphi(t)\leq \lim_{t\to\beta-0}\varphi(t)=+\infty,
\]
则我们有以下的换元公式:
\[
  \int_a^{+\infty}f(x)\mathrm dx
  =
  \int_\alpha^\beta f(\varphi(t))\varphi'(t)\mathrm dt.
\]

\noindent\textbf{注}\quad 上述等式是在以下意义下成立的: 当等式两端的积分有一个收敛时, 则另一个积分也收敛并且上述等式必成立; 若上面的两个积分中有一个发散, 则另外一个积分也发散. 以下若有两个无穷积分相等的等式, 则读者均可按上述意义来加以理解. 另外, 对 $\varphi(t)$ 是单调下降的情形, 请读者给出相应的换元公式。

现设函数 $u(x),v(x)$ 在 $[a,+\infty)$ 上连续可微, 且极限
\[
  u(+\infty)v(+\infty)\triangleq \lim_{x\to+\infty}u(x)v(x)
\]
存在, 则有以下分部积分公式
\[
  \int_a^{+\infty}u(x)v'(x)\mathrm dx
  =
  u(x)v(x)\big|_a^{+\infty}
  -
  \int_a^{+\infty}u'(x)v(x)\mathrm dx.
\]

读者可类似地给出无穷积分 $\int_{-\infty}^b f(x)\mathrm dx$ 及 $\int_{-\infty}^{+\infty}f(x)\mathrm dx$ 的换元公式与分部积分公式。

\noindent\textbf{例 8.1.4}\quad 计算无穷积分 $\int_1^{+\infty}\dfrac{\mathrm dx}{x\sqrt{1+x^2}}$。

\noindent\textbf{解}\quad 令 $x=\dfrac1t$, 则原积分可化为
\[
  \int_1^{+\infty}\frac{\mathrm dx}{x\sqrt{1+x^2}}
  =
  \int_0^1\frac{\mathrm dt}{\sqrt{1+t^2}}
  =
  \ln(t+\sqrt{1+t^2})\big|_0^1
  =
  \ln(1+\sqrt2).
\]
从上例可知, 一个无穷积分经过换元后有时可变成一个定积分。

\noindent\textbf{例 8.1.5}\quad 计算无穷积分 $\int_1^{+\infty}\dfrac{\arctan x}{x^3}\mathrm dx$。

\noindent\textbf{解}\quad 取 $u(x)=\arctan x,v(x)=-\dfrac1{2x^2}$, 则利用分部积分有
\[
\begin{aligned}
  \int_1^{+\infty}\frac{\arctan x}{x^3}\mathrm dx
  &=
  -\frac{\arctan x}{2x^2}\bigg|_1^{+\infty}
  +
  \int_1^{+\infty}\frac{\mathrm dx}{2x^2(1+x^2)}\\
  &=
  \frac{\pi}{8}
  +
  \frac12\int_1^{+\infty}\left(\frac1{x^2}-\frac1{1+x^2}\right)\mathrm dx\\
  &=
  \frac{\pi}{8}
  +
  \frac12\left(-\frac1x-\arctan x\right)\bigg|_1^{+\infty}
  =
  \frac12.
\end{aligned}
\]

\noindent\textbf{例 8.1.6}\quad 计算无穷积分
\[
  I_1=\int_0^{+\infty}e^{-ax}\sin bx\mathrm dx,\quad
  I_2=\int_0^{+\infty}e^{-ax}\cos bx\mathrm dx,\quad
  \text{其中 }a>0.
\]

\noindent\textbf{解}\quad 在下一节中, 我们将知道以上无穷积分是收敛的. 现在我们来求其值:
\[
\begin{aligned}
  I_1
  &=
  -\lim_{X\to+\infty}\frac1a\int_0^X\sin bx\,\mathrm d(e^{-ax})\\
  &=
  -\lim_{X\to+\infty}\frac1a\left(e^{-ax}\sin bx\big|_0^X
  -
  b\int_0^X e^{-ax}\cos bx\mathrm dx\right)\\
  &=
  \frac ba\int_0^{+\infty}e^{-ax}\cos bx\mathrm dx
  =
  \frac ba I_2,
\end{aligned}
\]
\[
\begin{aligned}
  I_2
  &=
  -\lim_{X\to+\infty}\frac1a\int_0^X\cos bx\,\mathrm d(e^{-ax})\\
  &=
  -\lim_{X\to+\infty}\frac1a\left(e^{-ax}\cos bx\big|_0^X
  +
  b\int_0^X e^{-ax}\sin bx\mathrm dx\right)\\
  &=
  \frac1a-\frac ba I_1.
\end{aligned}
\]
由以上两式可得
\[
  I_1=\frac b{a^2+b^2},\quad I_2=\frac a{a^2+b^2}.
\]

回忆一下, 我们说无穷积分 $\int_{-\infty}^{+\infty}f(x)\mathrm dx$ 收敛是指关于 $X_1$ 和 $X_2$ 的极限\allowbreak $\lim\limits_{\substack{X_1\to-\infty\\ X_2\to+\infty}}\int_{X_1}^{X_2}f(x)\mathrm dx$ 存在, 其中关于 $X_1$ 和 $X_2$ 的极限过程是两个独立的过程. 如果仅考虑 $\lim\limits_{X\to+\infty}\int_{-X}^X f(x)\mathrm dx$, 则当它收敛时, 我们称 $\int_{-\infty}^{+\infty}f(x)\mathrm dx$ 在柯西主值意义下是收敛的, 此时我们用
\[
  \mathrm{V.P.}\int_{-\infty}^{+\infty}f(x)\mathrm dx
\]
来表示极限 $\lim\limits_{X\to+\infty}\int_{-X}^X f(x)\mathrm dx$, 同时称该极限值为积分 $\int_{-\infty}^{+\infty}f(x)\mathrm dx$ 的柯西主值. 显然, 对任何 $(-\infty,+\infty)$ 上可积的奇函数 $f(x)$, 有
\[
  \mathrm{V.P.}\int_{-\infty}^{+\infty}f(x)\mathrm dx=0.
\]

从某种意义上来说, 当讨论无穷积分 $\int_{-\infty}^{+\infty}f(x)\mathrm dx$ 的敛散性问题时, 我们必须讨论两个独立的极限过程. 而讨论该无穷积分的柯西主值时, 我们只要讨论一个极限过程. 无穷积分的柯西主值在一些数学分支中有应用。

\section{无穷积分敛散性的判别法}

在对无穷积分 $\int_a^{+\infty}f(x)\mathrm dx$ 的讨论中, 我们假定了函数 $f(x)$ 在 $[a,+\infty)$ 的有限子区间上都是可积的, 而要研究的 $\int_a^{+\infty}f(x)\mathrm dx$ 敛散性问题, 实际上就是研究当 $X\to+\infty$ 时连续函数 $F(X)=\int_a^X f(x)\mathrm dx$ 的极限是否存在的问题. 因此从理论上讲, 无穷积分的敛散性可以认为是已经解决的问题. 但是人们很快就发现, 对于许多的函数 $f(x)$, 相应的 $F(X)$ 并不能容易求出, 因此我们要寻找直接从被积函数的性质来判定广义积分是否收敛的方法。

下面我们主要讨论 $\int_a^{+\infty}f(x)\mathrm dx$ 的敛散性问题, 关于 $\int_{-\infty}^b f(x)\mathrm dx$ 及 $\int_{-\infty}^{+\infty}f(x)\mathrm dx$ 的敛散性, 请读者自己给出相应的结论。

首先由无穷积分 $\int_a^{+\infty}f(x)\mathrm dx$ 收敛的定义, 我们容易得到下面的柯西准则。

\noindent\textbf{定理 8.2.1（柯西准则）}\quad 设函数 $f(x)$ 在 $[a,+\infty)$ 上有定义, 对于 $\forall X>a$, 在区间 $[a,X]$ 上可积, 则无穷积分 $\int_a^{+\infty}f(x)\mathrm dx$ 收敛的充分必要条件是: 对于 $\forall\varepsilon>0$, $\exists M>a$, 当 $X''>X'>M$ 时, 有
\[
  \left|\int_{X'}^{X''}f(x)\mathrm dx\right|<\varepsilon.
\]

此定理的证明可以从无穷积分敛散性的定义直接推出, 故省略。

\noindent\textbf{定义 8.2.1}\quad 设函数 $f(x)$ 在 $[a,+\infty)$ 上有定义, 对于 $\forall X>a$, 在区间 $[a,X]$ 上可积. 若无穷积分 $\int_a^{+\infty}|f(x)|\mathrm dx$ 收敛, 则称无穷积分 $\int_a^{+\infty}f(x)\mathrm dx$ 绝对收敛; 若无穷积分 $\int_a^{+\infty}f(x)\mathrm dx$ 收敛, 但无穷积分 $\int_a^{+\infty}|f(x)|\mathrm dx$ 发散, 则称无穷积分 $\int_a^{+\infty}f(x)\mathrm dx$ 条件收敛。

\noindent\textbf{定理 8.2.2}\quad 设函数 $f(x)$ 在 $[a,+\infty)$ 上有定义, 对于 $\forall X>a$, 在区间 $[a,X]$ 上可积. 若无穷积分 $\int_a^{+\infty}f(x)\mathrm dx$ 绝对收敛, 则无穷积分 $\int_a^{+\infty}f(x)\mathrm dx$ 必收敛。

\noindent\textbf{证明}\quad 假设 $\int_a^{+\infty}|f(x)|\mathrm dx$ 收敛, 由柯西准则, 对于 $\forall\varepsilon>0$, $\exists M>a$, 当 $X''>X'>M$ 时, 有
\[
  \left|\int_{X'}^{X''}|f(x)|\mathrm dx\right|
  =
  \int_{X'}^{X''}|f(x)|\mathrm dx<\varepsilon.
\]
因此, 当 $X''>X'>M$ 时, 有
\[
  \left|\int_{X'}^{X''}f(x)\mathrm dx\right|
  \leq
  \int_{X'}^{X''}|f(x)|\mathrm dx<\varepsilon.
\]
再由柯西准则知 $\int_a^{+\infty}f(x)\mathrm dx$ 收敛. 证毕。

从上述定理知, 无穷积分是否绝对收敛的问题可以转化为非负函数的无穷积分的敛散性问题。

\noindent\textbf{定理 8.2.3}\quad 设非负函数 $f(x)$ 在 $[a,+\infty)$ 上有定义, 对于 $\forall X>a$, 在 $[a,X]$ 上可积, 则无穷积分 $\int_a^{+\infty}f(x)\mathrm dx$ 收敛的充分必要条件是: 存在 $A>0$, 使得对一切 $X\geq a$, 有
\[
  \int_a^X f(x)\mathrm dx\leq A.
\]

\noindent\textbf{证明}\quad 事实上, 由定积分关于积分区间的可加性及 $f(x)\geq0$ 知, $F(X)=\int_a^X f(x)\mathrm dx$ 在 $[a,+\infty)$ 上是 $X$ 的单调上升函数. 因此由单调有界定理即可推出定理 8.2.3. 证毕。

对于非负函数的无穷积分, 我们有下面的比较定理。

\noindent\textbf{定理 8.2.4}\quad 设非负函数 $f(x),g(x)$ 在 $[a,+\infty)$ 上有定义, 且对于 $\forall X>a$, 在 $[a,X]$ 上可积. 若存在常数 $c_1>0,c_2>0$ 及 $M_0\geq a$, 使得当 $x\geq M_0$ 时, 成立不等式
\[
  c_1f(x)\leq c_2g(x),
\]
则我们有下述结论:

(1) 若 $\int_a^{+\infty}g(x)\mathrm dx$ 收敛, 则 $\int_a^{+\infty}f(x)\mathrm dx$ 也收敛;

(2) 若 $\int_a^{+\infty}f(x)\mathrm dx$ 发散, 则 $\int_a^{+\infty}g(x)\mathrm dx$ 也发散。

\noindent\textbf{证明}\quad 我们只证 (1), (2) 的证明留给读者。

由假设 $\int_a^{+\infty}g(x)\mathrm dx$ 收敛, 根据柯西准则, 对于 $\forall\varepsilon>0$, $\exists M_1>0$, 当 $X''>X'>M_1$ 时, 有
\[
  \int_{X'}^{X''}g(x)\mathrm dx<\frac{c_1}{c_2}\varepsilon.
\]
由于当 $x\geq M_0$ 时, 有
\[
  f(x)\leq \frac{c_2}{c_1}g(x),
\]
因此我们取 $M=\max\{M_0,M_1\}$, 当 $X''>X'>M$ 时, 有
\[
  0\leq \int_{X'}^{X''}f(x)\mathrm dx
  \leq
  \frac{c_2}{c_1}\int_{X'}^{X''}g(x)\mathrm dx<\varepsilon.
\]
(1) 得证. 证毕。

\noindent\textbf{推论}\quad 设非负函数 $f(x),g(x)(g(x)\ne0)$ 在 $[a,+\infty)$ 上有定义, 且对于 $\forall X>a$, 在区间 $[a,X]$ 上可积. 若
\[
  \lim_{x\to+\infty}\frac{f(x)}{g(x)}=l,
\]
则

(1) 当 $0<l<+\infty$ 时, $\int_a^{+\infty}g(x)\mathrm dx$ 与 $\int_a^{+\infty}f(x)\mathrm dx$ 同时收敛或同时发散;

(2) 当 $l=0$ 时, 若 $\int_a^{+\infty}g(x)\mathrm dx$ 收敛, 则 $\int_a^{+\infty}f(x)\mathrm dx$ 收敛;

(3) 当 $l=+\infty$ 时, 若 $\int_a^{+\infty}g(x)\mathrm dx$ 发散, 则 $\int_a^{+\infty}f(x)\mathrm dx$ 发散。

上述推论可由极限的保号性和定理 8.2.4 推出。

显然, 对两个非负函数的无穷积分, 以上定理和推论告诉我们, 取值大的函数的无穷积分收敛, 则取值小的函数的无穷积分也收敛; 反之, 取值小的函数的无穷积分发散, 则取值大的函数的无穷积分也发散。

在判断一些无穷积分的敛散性时, 我们有一些 ``标准函数'', 其他的一些函数可与它们作比较. 这些 ``标准函数'' 有 $\dfrac1{x^p},\dfrac1{x^p\ln^q x},e^{-ax}$ 等. 读者应该熟悉这些函数中参数取哪些值是使得无穷积分收敛的, 而哪些值是使得无穷积分发散的. 例如, 在 $\int_2^{+\infty}\dfrac{\mathrm dx}{x^p\ln^q x}$ 中, 当 $p>1$ 时, 对任意的 $q\in\mathbb R$, 无穷积分收敛; 而当 $p=1$ 时, 对于 $q>1$ 收敛, 对于其他的情形该无穷积分均发散. 再如, 在 $\int_0^{+\infty}P(x)e^{-ax}\mathrm dx$ 中(其中 $P(x)$ 是一个多项式), 当 $a>0$ 时, 无穷积分是收敛的。

在上节例 8.1.6 中, 由于被积函数在 $[0,+\infty)$ 上满足 $|e^{-ax}\sin bx|\leq e^{-ax}$ 和 $|e^{-ax}\cos bx|\leq e^{-ax}$, 因此 $I_1=\int_0^{+\infty}e^{-ax}\sin bx\mathrm dx$ 与 $I_2=\int_0^{+\infty}e^{-ax}\cos bx\mathrm dx$ 均是绝对收敛的。

\noindent\textbf{例 8.2.1}\quad 讨论无穷积分 $\int_0^{+\infty}\dfrac{x^2}{x^n+x^2+1}\mathrm dx$ 的敛散性, 其中 $n$ 为一正整数。

\noindent\textbf{解}\quad 由于
\[
  \lim_{x\to+\infty}
  \frac{\dfrac{x^2}{x^n+x^2+1}}{\dfrac1{x^{n-2}}}
  =
  1,
\]
而当 $n-2>1$, 即 $n>3$ 时, $\int_0^{+\infty}\dfrac1{x^{n-2}}\mathrm dx$ 收敛, 因此当 $n>3$ 时, 原无穷积分收敛; 而当 $n\leq3$ 时, 无穷积分发散。

若一个有理函数 $R(x)=\dfrac{P(x)}{Q(x)}$ 在 $[a,+\infty)$ 上有定义, 其中 $P(x)$ 与 $Q(x)$ 为两个既约多项式, 它们的次数分别为 $m,n$, 则 $\int_a^{+\infty}R(x)\mathrm dx$ 收敛的充分必要条件是 $n-m\geq2$. 对于有理函数的无穷积分, 虽然我们很容易判断其收敛性, 但计算其值并不容易. 在复变函数课程中我们有系统的办法来求其值。

\noindent\textbf{例 8.2.2}\quad 讨论无穷积分 $\int_1^{+\infty}\dfrac{\cos x\sin\dfrac1x}{x}\mathrm dx$ 的敛散性。

\noindent\textbf{解}\quad 由于
\[
  \left|\frac{\cos x\sin\dfrac1x}{x}\right|
  \leq
  \frac1x\sin\frac1x,\quad x\geq1,
\]
且
\[
  \lim_{x\to+\infty}
  \frac{\dfrac1x\sin\dfrac1x}{\dfrac1{x^2}}=1,
\]
而 $\int_1^{+\infty}\dfrac1{x^2}\mathrm dx$ 收敛, 所以 $\int_1^{+\infty}\dfrac{\cos x\sin\dfrac1x}{x}\mathrm dx$ 绝对收敛。

\noindent\textbf{例 8.2.3}\quad 讨论无穷积分 $\int_1^{+\infty}\left[\dfrac1{x^p}-\ln\left(1+\dfrac1{x^p}\right)\right]\mathrm dx\ (p>0)$ 的敛散性。

\noindent\textbf{解}\quad 在 $\ln(1+t)$ 的麦克劳林展开式中令 $t=\dfrac1{x^p}$, 我们有
\[
  \ln\left(1+\frac1{x^p}\right)
  =
  \frac1{x^p}-\frac1{2x^{2p}}+o\left(\frac1{x^{2p}}\right)
  \quad(x\to+\infty),
\]
从而有
\[
  0\leq
  \frac1{x^p}-\ln\left(1+\frac1{x^p}\right)
  =
  \frac1{2x^{2p}}+o\left(\frac1{x^{2p}}\right)
  \sim
  \frac1{2x^{2p}}
  \quad(x\to+\infty).
\]
因此, $\int_1^{+\infty}\left[\dfrac1{x^p}-\ln\left(1+\dfrac1{x^p}\right)\right]\mathrm dx$ 当 $p>\dfrac12$ 时收敛, 当 $p\leq\dfrac12$ 时发散。

以上我们讨论了非负函数无穷积分的敛散性问题. 正如我们所指出的那样, 这些判别法则可以用来判定无穷积分是否绝对收敛. 值得注意的是, 当 $f(x)\geq0$ 时, 由 $\int_a^{+\infty}f(x)\mathrm dx$ 的几何意义来看, 似乎当 $x$ 越大时, $f(x)$ 的值要很小时才能保证无穷积分 $\int_a^{+\infty}f(x)\mathrm dx$ 的收敛性. 当 $f(x)$ 不连续时, 我们容易看出这个想法是不对的, 这可以从以下事实得到验证. 取严格单调上升序列 $\{x_n\}\subset[a,+\infty)$, 使得 $\lim\limits_{n\to\infty}x_n=+\infty$, 由于改变 $f(x)$ 在 $\{x_n\}$ 上的取值并不改变 $f(x)$ 在任何区间的可积性, 因此我们可以使得 $f(x)$ 在该序列的值改变后的函数在这一序列上的值趋于 $\infty$. 即使 $f(x)$ 是连续的, 我们也容易构造出一个在 $[a,+\infty)$ 上无界的函数 $f(x)$, 使得广义积分 $\int_a^{+\infty}f(x)\mathrm dx$ 收敛。

\noindent\textbf{例 8.2.4}\quad 设函数 $f(x)$ 在 $[1,+\infty)$ 上的定义如下:
\[
  f(x)=
  \begin{cases}
    \dfrac1{x^2}, & x\in[1,x_2'],\\
    \dfrac1{x^2}, & x\in[x_n',x_{n+1}],\\
    f(x_{n+1})+3^{n+1}\bigl[2^{n+1}-f(x_{n+1})\bigr](x-x_{n+1}), & x\in(x_{n+1},n+1],\\
    2^{n+1}-3^{n+1}\bigl[2^{n+1}-f(x_{n+1}')\bigr](x-n-1), & x\in(n+1,x_{n+1}'],
  \end{cases}
\]
其中 $x_n=n-\dfrac1{3^n},x_n'=n+\dfrac1{3^n},n=2,3,\cdots$ ($f(x)$ 的大致图像如图 8.2.1 所示), 证明无穷积分 $\int_1^{+\infty}f(x)\mathrm dx$ 收敛。

\noindent\textbf{证明}\quad 从 $f(x)$ 的构造可以看出 $f(x)>0(x\in[1,+\infty)), f(x)\in C[1,+\infty)$, 并且 $\lim\limits_{n\to\infty}f(n)=+\infty$. 对于 $\forall X>1$, 取正整数 $N>X$, 则有
\[
\begin{aligned}
  \int_1^X f(x)\mathrm dx
  &\leq
  \int_1^N f(x)\mathrm dx
  \leq
  \int_1^N\frac{\mathrm dx}{x^2}
  +
  \sum_{n=2}^N\int_{x_n}^{x_n'}f(x)\mathrm dx\\
  &\leq
  1+\sum_{n=2}^N\int_{x_n}^{x_n'}2^n\mathrm dx
  <
  1+2\sum_{n=2}^N\frac{2^n}{3^n}<4.
\end{aligned}
\]
因此无穷积分 $\int_1^{+\infty}f(x)\mathrm dx$ 收敛。

作为练习, 读者可以证明: 当 $f(x)$ 在 $[a,+\infty)$ 上单调时, 若无穷积分 $\int_a^{+\infty}f(x)\mathrm dx$ 收敛, 则必有 $\lim\limits_{x\to+\infty}f(x)=0$。

现在我们转向讨论条件收敛的无穷积分. 对于这类无穷积分, 被积函数 $f(x)$ 的绝对值的无穷积分趋于 $+\infty$. 若 $f(x)$ 在 $[a,+\infty)$ 上连续, 当 $x\to+\infty$ 时, $f(x)$ 取值必须在 $x$ 轴的两侧波动, 使得其图像与 $x$ 轴所围的图形在 $x$ 轴的上面与下面的面积不断抵消, 才能使得 $\int_a^{+\infty}f(x)\mathrm dx$ 收敛. 在此方面我们有以下两个有用的判别法。

\noindent\textbf{定理 8.2.5（狄利克雷判别法）}\quad 设函数 $f(x),g(x)$ 在 $[a,+\infty)$ 上有定义, 且满足下面两个条件:

(1) 对于 $\forall X>a$, $f(x)$ 在区间 $[a,X]$ 上可积, 并且 $\exists M>0$, 使得对于 $\forall X>a$, 有
\[
  \left|\int_a^X f(x)\mathrm dx\right|\leq M;
\]

(2) $g(x)$ 在 $[a,+\infty)$ 单调, 并且 $\lim\limits_{x\to+\infty}g(x)=0$,

则无穷积分 $\int_a^{+\infty}f(x)g(x)\mathrm dx$ 收敛。

\noindent\textbf{证明}\quad 对于 $\forall\varepsilon>0$, 由 $\lim\limits_{x\to+\infty}g(x)=0$, $\exists X>0$, 当 $x>X$ 时, 有
\[
  |g(x)|<\frac{\varepsilon}{4M}.
\]
对于任意的 $X''>X'>X$, 由条件 (1) 和 (2), 我们可以对定积分 $\int_{X'}^{X''}f(x)g(x)\mathrm dx$ 应用定积分第二中值定理, 从而在 $(X',X'')$ 内存在 $\xi$, 使得
\[
\begin{aligned}
  \left|\int_{X'}^{X''}f(x)g(x)\mathrm dx\right|
  &=
  \left|g(X')\int_{X'}^\xi f(x)\mathrm dx
  +
  g(X'')\int_\xi^{X''}f(x)\mathrm dx\right|\\
  &<
  \frac{\varepsilon}{4M}
  \left\{
  \left|\int_a^\xi f(x)\mathrm dx-\int_a^{X'}f(x)\mathrm dx\right|
  +
  \left|\int_a^{X''}f(x)\mathrm dx-\int_a^\xi f(x)\mathrm dx\right|
  \right\}\\
  &\leq \frac{\varepsilon}{4M}\cdot4M=\varepsilon.
\end{aligned}
\]
由柯西准则我们知无穷积分 $\int_a^{+\infty}f(x)g(x)\mathrm dx$ 收敛. 证毕。

\noindent\textbf{定理 8.2.6（阿贝尔判别法）}\quad 设函数 $f(x),g(x)$ 在 $[a,+\infty)$ 上有定义, 并且满足下面两个条件:

(1) 对于 $\forall X>a$, $f(x)$ 在 $[a,X]$ 上可积, 并且 $\int_a^{+\infty}f(x)\mathrm dx$ 收敛;

(2) $g(x)$ 在 $[a,+\infty)$ 单调有界,

则无穷积分 $\int_a^{+\infty}f(x)g(x)\mathrm dx$ 收敛。

\noindent\textbf{证明}\quad 由条件 (2) 我们知 $\lim\limits_{x\to+\infty}g(x)$ 存在, 设该极限值为 $B$. 容易验证 $f(x)$ 和 $G(x)=g(x)-B$ 满足狄利克雷判别法(定理 8.2.5) 的所有条件, 从而 $\int_a^{+\infty}f(x)G(x)\mathrm dx$ 收敛. 由
\[
  \int_a^{+\infty}f(x)g(x)\mathrm dx
  =
  \int_a^{+\infty}f(x)G(x)\mathrm dx
  +
  B\int_a^{+\infty}f(x)\mathrm dx
\]
即知定理结论成立. 证毕。

\noindent\textbf{例 8.2.5}\quad 证明无穷积分 $\int_{-\infty}^{+\infty}\dfrac{\sin x}{x}\mathrm dx$ 条件收敛。

\noindent\textbf{解}\quad 注意到被积函数为偶函数以及
\[
  f(x)=
  \begin{cases}
    \dfrac{\sin x}{x}, & x\ne0,\\
    1, & x=0
  \end{cases}
\]
在 $(-\infty,+\infty)$ 上连续, 因此我们只要证明无穷积分 $\int_{2\pi}^{+\infty}\dfrac{\sin x}{x}\mathrm dx$ 条件收敛即可。

对于 $\forall X>2\pi$, 我们有
\[
  \left|\int_{2\pi}^X\sin x\mathrm dx\right|
  =
  |\cos2\pi-\cos X|<2,
\]
又由于 $\dfrac1x$ 在 $[2\pi,+\infty)$ 上单调下降且 $\lim\limits_{x\to+\infty}\dfrac1x=0$, 由狄利克雷判别法知 $\int_{2\pi}^{+\infty}\dfrac{\sin x}{x}\mathrm dx$ 收敛。

由于
\[
  \frac{|\sin x|}{x}\geq \frac{\sin^2x}{x}=\frac1{2x}-\frac{\cos x}{2x},
\]
类似上面讨论, 我们可证 $\int_{2\pi}^{+\infty}\dfrac{\cos x}{2x}\mathrm dx$ 收敛. 由于 $\int_{2\pi}^{+\infty}\dfrac{\mathrm dx}{2x}$ 发散, 因此 $\int_{2\pi}^{+\infty}\dfrac{\sin^2x}{x}\mathrm dx$ 发散. 再由比较判别法知 $\int_{2\pi}^{+\infty}\dfrac{|\sin x|}{x}\mathrm dx$ 发散. 因此无穷积分 $\int_{2\pi}^{+\infty}\dfrac{\sin x}{x}\mathrm dx$ 条件收敛。

从此例子可以清楚地看出, 当 $x$ 充分大时, 由 $\sin x$ 的周期性, 对曲线 $y=\dfrac{\sin x}{x}$ 与 $x$ 轴所围的每一小块的面积 $A_i(i=1,2,\cdots)$(见图 8.2.2), 有
\[
  A_1>A_2>A_3>\cdots,
\]
且 $\lim\limits_{n\to+\infty}A_n=0$. 由于 $x$ 轴上面小块与下面小块的面积不断抵消, 所以 $\int_{2\pi}^{+\infty}\dfrac{\sin x}{x}\mathrm dx$ 收敛。

另外, 在不定积分中我们知道 $\int\dfrac{\sin x}{x}\mathrm dx$ 不能用初等函数表示, 因此对 $-\infty<a<b<+\infty$, 定积分 $\int_a^b\dfrac{\sin x}{x}\mathrm dx$ 就很难计算出来. 但以后我们可以用多种方法求出 $\int_0^{+\infty}\dfrac{\sin x}{x}\mathrm dx$ 的值. 类似的例子还有 $\int_0^{+\infty}e^{-x^2}\mathrm dx$ 等。

\noindent\textbf{例 8.2.6}\quad 试讨论 $I=\int_1^{+\infty}\ln\left(1+\dfrac{\sin x}{x^p}\right)\mathrm dx(p>0)$ 的敛散性。

\noindent\textbf{解}\quad 由
\[
  \ln(1+t)=t-\frac{t^2}{2}+o(t^2)\quad(t\to0)
\]
知
\[
\begin{aligned}
  \ln\left(1+\frac{\sin x}{x^p}\right)
  &=
  \frac{\sin x}{x^p}
  -
  \frac{\sin^2x}{2x^{2p}}
  +
  o\left(\frac1{x^{2p}}\right)\\
  &=
  \frac{\sin x}{x^p}
  +
  \frac{\cos2x}{4x^{2p}}
  -
  [1-o(1)]\frac1{2x^{2p}}
  \quad(x\to+\infty).
\end{aligned}
\]
因此我们有
\[
\begin{aligned}
  I
  &=
  \int_1^{+\infty}\ln\left(1+\frac{\sin x}{x^p}\right)\mathrm dx\\
  &=
  \int_1^{+\infty}\frac{\sin x}{x^p}\mathrm dx
  +
  \int_1^{+\infty}\frac{\cos2x}{4x^{2p}}\mathrm dx
  -
  \int_1^{+\infty}[1-o(1)]\frac{\mathrm dx}{2x^{2p}}\\
  &\triangleq I_1+I_2-I_3.
\end{aligned}
\]
当 $0<p\leq\dfrac12$ 时, 由狄利克雷判别法知 $I_1$ 和 $I_2$ 收敛, 而注意到此时 $I_3$ 发散, 因此 $I$ 发散; 当 $1\geq p>\dfrac12$ 时, $I_1$ 条件收敛, 而 $I_2$ 和 $I_3$ 绝对收敛, 因此 $I$ 条件收敛; 当 $p>1$ 时, $I_1,I_2$ 和 $I_3$ 均绝对收敛, 从而 $I$ 绝对收敛。

\noindent\textbf{注}\quad 注意到对所有的 $p>0$, 我们都有
\[
  \lim_{x\to+\infty}
  \frac{\ln\left(1+\dfrac{\sin x}{x^p}\right)}{\dfrac{\sin x}{x^p}}
  =
  1.
\]
此例说明若被积函数不保号时, 不能用比较判别法来确定它是否收敛。

\section{瑕积分}

\subsection{瑕积分的概念}

正如我们在本章开始所指出的, 定积分的第二个推广是研究在有限区间上无界函数的积分问题. 回忆一下, 在研究函数可积性时, 我们知道
\[
  F(x)=
  \begin{cases}
    x^2\sin\dfrac1{x^2}, & x\ne0,\\
    0, & x=0
  \end{cases}
\]
是函数
\[
  f(x)=
  \begin{cases}
    2x\sin\dfrac1{x^2}-\dfrac2x\cos\dfrac1{x^2}, & x\ne0,\\
    0, & x=0
  \end{cases}
\]
的一个原函数. 函数 $f(x)$ 是无界的, 因此它在包含 $x=0$ 的任何闭区间 $[a,b]$ 上不可积, 但此时函数 $f(x)$ 却有一个有界的原函数. 在本节中, 我们希望研究此类函数的积分问题. 我们称这类函数的积分为瑕积分. 因为 $f(x)$ 仅仅在 $x=0$ 附近无界, 所以 $x=0$ 也称为 $f(x)$ 的一个瑕点. 下面我们称 $x_0$ 是 $f(x)$ 的一个瑕点即是指 $f(x)$ 在 $x_0$ 的某个去心(左或右)邻域内有定义, 但在该去心(左或右)邻域内无界。

首先我们讨论函数 $f(x)$ 在区间 $[a,b]$ 上只有一个瑕点的情形。

\noindent\textbf{定义 8.3.1}\quad 设函数 $f(x)$ 在区间 $(a,b]$ 上有定义, $a$ 是 $f(x)$ 的一个瑕点. 若对于 $\forall 0<\delta<b-a$, $f(x)$ 在区间 $[a+\delta,b]$ 上可积, 且极限
\begin{equation}
  \lim_{\delta\to+0}\int_{a+\delta}^b f(x)\mathrm dx
  \tag{8.3.1}
\end{equation}
存在, 则称瑕积分 $\int_a^b f(x)\mathrm dx$ 收敛, 并记
\[
  \int_a^b f(x)\mathrm dx
  =
  \lim_{\delta\to+0}\int_{a+\delta}^b f(x)\mathrm dx;
\]

若极限 (8.3.1) 不存在, 则称瑕积分 $\int_a^b f(x)\mathrm dx$ 发散。

如果 $b$ 为函数 $f(x)$ 的瑕点, 我们可以类似定义
\[
  \int_a^b f(x)\mathrm dx
  =
  \lim_{\delta\to+0}\int_a^{b-\delta}f(x)\mathrm dx.
\]

当 $c\in(a,b)$ 为 $f(x)$ 在 $[a,b]$ 上的唯一瑕点时, 我们称 $\int_a^b f(x)\mathrm dx$ 收敛是指瑕积分 $\int_a^c f(x)\mathrm dx$ 与 $\int_c^b f(x)\mathrm dx$ 同时收敛. 此时, 我们还可类似地定义柯西主值意义下的收敛(请读者给出)。

当 $a$ 是函数 $f(x)$ 在区间 $[a,b]$ 上的唯一瑕点时, 若 $F(x)$ 是 $f(x)$ 在 $(a,b]$ 上的一个原函数, 则瑕积分 $\int_a^b f(x)\mathrm dx$ 可由下式计算出:
\[
  \int_a^b f(x)\mathrm dx
  =
  \lim_{\delta\to+0}\int_{a+\delta}^b f(x)\mathrm dx
  =
  F(b)-\lim_{\delta\to+0}F(a+\delta).
\]
例如, 对于前面的例子我们有
\[
  \int_0^1\left(2x\sin\frac1{x^2}-\frac2x\cos\frac1{x^2}\right)\mathrm dx
  =
  \lim_{\delta\to+0}x^2\sin\frac1{x^2}\bigg|_\delta^1
  =
  \sin1.
\]
以上的分析实际上给出了我们计算瑕积分的一个方法. 对于瑕积分, 我们也有相应的换元公式与分部积分公式, 这些公式请读者自己给出。

\noindent\textbf{例 8.3.1}\quad 计算瑕积分 $\int_{-1}^1\dfrac{\arccos x}{\sqrt{1-x^2}}\mathrm dx$。

\noindent\textbf{解}\quad 由洛必达法则我们有
\[
  \lim_{x\to1-0}\frac{\arccos x}{\sqrt{1-x^2}}
  =
  \lim_{x\to1-0}
  \frac{-\dfrac1{\sqrt{1-x^2}}}{-\dfrac{2x}{2\sqrt{1-x^2}}}
  =
  1,
\]
因此被积函数有唯一的瑕点 $x=-1$. 所以有
\[
  \int_{-1}^1\frac{\arccos x}{\sqrt{1-x^2}}\mathrm dx
  =
  \lim_{\delta\to+0}\left(-\frac12\arccos^2x\right)\bigg|_{-1+\delta}^1
  =
  \frac{\pi^2}{2}.
\]

\noindent\textbf{例 8.3.2}\quad 计算瑕积分 $\int_0^1\ln x\mathrm dx$。

\noindent\textbf{解}\quad 容易看出 $x=0$ 是瑕点. 利用分部积分法得
\[
  \int_0^1\ln x\mathrm dx
  =
  \lim_{\delta\to+0}x\ln x\big|_\delta^1
  -
  \int_0^1\mathrm dx
  =
  -1.
\]

\noindent\textbf{例 8.3.3}\quad 讨论瑕积分 $\int_a^b\dfrac{\mathrm dx}{(b-x)^p}$ 的敛散性。

\noindent\textbf{解法 1}\quad 当 $p=1$ 时,
\[
\begin{aligned}
  \int_a^b\frac{\mathrm dx}{b-x}
  &=
  \lim_{\delta\to+0}\int_a^{b-\delta}\frac{\mathrm dx}{b-x}
  =
  \lim_{\delta\to+0}[-\ln(b-x)]\big|_a^{b-\delta}\\
  &=
  \lim_{\delta\to+0}[\ln(b-a)-\ln\delta]=+\infty.
\end{aligned}
\]
当 $p\ne1$ 时,
\[
\begin{aligned}
  \int_a^b\frac{\mathrm dx}{(b-x)^p}
  &=
  \lim_{\delta\to+0}\int_a^{b-\delta}\frac{\mathrm dx}{(b-x)^p}\\
  &=
  -\lim_{\delta\to+0}\frac1{1-p}(b-x)^{1-p}\bigg|_a^{b-\delta}\\
  &=
  \begin{cases}
    \dfrac{(b-a)^{1-p}}{1-p}, & p<1,\\
    +\infty, & p>1.
  \end{cases}
\end{aligned}
\]
因此, $\int_a^b\dfrac{\mathrm dx}{(b-x)^p}$ 当 $p<1$ 时收敛, 当 $p\geq1$ 时发散。

\noindent\textbf{解法 2}\quad 我们也可以用以下办法来讨论 $\int_a^b\dfrac{\mathrm dx}{(b-x)^p}$ 的敛散性. 作变换 $t=\dfrac1{b-x}$, 则当 $x=a$ 时, 有 $t=\dfrac1{b-a}$, 当 $x\to b-0$ 时, 有 $t\to+\infty$, 并且
\[
  \mathrm dt=\frac{\mathrm dx}{(b-x)^2}=t^2\mathrm dx,\quad
  \mathrm dx=t^{-2}\mathrm dt.
\]
将它们代入 $\int_a^b\dfrac{\mathrm dx}{(b-x)^p}$, 得
\[
  \int_a^b\frac{\mathrm dx}{(b-x)^p}
  =
  \int_{\frac1{b-a}}^{+\infty}t^{-2}t^p\mathrm dt
  =
  \int_{\frac1{b-a}}^{+\infty}\frac{\mathrm dt}{t^{2-p}}.
\]
由无穷积分的结果知, 当 $2-p>1$, 即 $p<1$ 时, $\int_a^b\dfrac{\mathrm dx}{(b-x)^p}$ 收敛; 当 $2-p\leq1$, 即 $p\geq1$ 时, $\int_a^b\dfrac{\mathrm dx}{(b-x)^p}$ 发散。

一般来说, 通过类似于上述解法 2 中的变换, 一个瑕积分总可以转化成一个无穷积分. 另外, 一个无穷积分 $\int_a^{+\infty}f(x)\mathrm dx(a>0)$ 通过变换 $x=\dfrac1t$ 可以化为 $\int_0^{\frac1a}\dfrac1{t^2}f\left(\dfrac1t\right)\mathrm dt$, 因此一个无穷积分就有可能化成一个定积分或者瑕积分。

\subsection{瑕积分敛散性的判别法}

对于瑕积分的敛散性的判定, 一个方法是将其化为无穷积分来讨论. 但有时由于化成无穷积分后被积函数会更加复杂, 因此我们需要讨论直接从瑕积分本身来判定的方法. 以下我们列出函数 $f(x)$ 在区间 $[a,b]$ 中仅有瑕点 $b$ 时瑕积分的有关判定定理, 其证明则留给读者. 因此, 在本小节中, 定义在 $[a,b)$ 上的函数 $f(x),g(x)$ 等, 总假定 $b$ 是它们的瑕点, 并且它们在 $[a,b)$ 的任何子区间上可积. 另外, 我们要指出的是, 对于瑕积分, 同样有绝对收敛、条件收敛等概念和相应的结论。

\noindent\textbf{定理 8.3.1（柯西准则）}\quad 瑕积分 $\int_a^b f(x)\mathrm dx$ 收敛的充分必要条件是: 对于 $\forall\varepsilon>0$, $\exists\delta>0$, 当 $0<\delta''<\delta'<\delta$ 时, 有
\[
  \left|\int_{b-\delta'}^{b-\delta''}f(x)\mathrm dx\right|<\varepsilon.
\]

\noindent\textbf{定理 8.3.2}\quad 设非负函数 $f(x),g(x)$ 在区间 $[a,b)$ 上满足: 存在正常数 $c_1,c_2$, 使得当 $x\in[b-\delta_0,b)(0<\delta_0<b-a)$ 时, 有
\[
  c_1f(x)\leq c_2g(x),
\]
则

(1) 若 $\int_a^b g(x)\mathrm dx$ 收敛, 则 $\int_a^b f(x)\mathrm dx$ 也收敛;

(2) 若 $\int_a^b f(x)\mathrm dx$ 发散, 则 $\int_a^b g(x)\mathrm dx$ 也发散。

\noindent\textbf{推论}\quad 设非负函数 $f(x),g(x)(g(x)\ne0)$ 满足
\[
  \lim_{x\to b-0}\frac{f(x)}{g(x)}=l,
\]
则

(1) 当 $0<l<+\infty$ 时, $\int_a^b f(x)\mathrm dx$ 与 $\int_a^b g(x)\mathrm dx$ 同时收敛或同时发散;

(2) 当 $l=0$ 时, 若 $\int_a^b g(x)\mathrm dx$ 收敛, 则 $\int_a^b f(x)\mathrm dx$ 收敛;

(3) 当 $l=+\infty$ 时, 若 $\int_a^b g(x)\mathrm dx$ 发散, 则 $\int_a^b f(x)\mathrm dx$ 发散。

同样, 对于瑕积分, 我们也有一些 ``标准函数'' 供其他函数来比较, 如 $\dfrac1{(b-x)^p}$。

\noindent\textbf{例 8.3.4}\quad 讨论瑕积分 $\int_{-1}^1\dfrac{\mathrm dx}{(1-x^2)^p}$ 的敛散性。

\noindent\textbf{解}\quad 非负被积函数有两个瑕点 $x=\pm1$. 由于
\[
  \lim_{x\to1-0}
  \frac{\dfrac1{(1-x^2)^p}}{\dfrac1{(1-x)^p}}
  =
  \frac1{2^p},\quad
  \lim_{x\to-1+0}
  \frac{\dfrac1{(1-x^2)^p}}{\dfrac1{(1+x)^p}}
  =
  \frac1{2^p},
\]
因此可知, 当 $p<1$ 时, 瑕积分收敛; 当 $p\geq1$ 时, 瑕积分发散。

对于瑕积分, 我们还有以下的判别法。

\noindent\textbf{定理 8.3.3（狄利克雷判别法）}\quad 设函数 $f(x),g(x)$ 在区间 $[a,b)$ 上满足下述条件:

(1) $\exists M>0$, 使得对于 $\forall\delta>0$, 有
\[
  \left|\int_a^{b-\delta}f(x)\mathrm dx\right|\leq M;
\]

(2) $g(x)$ 在 $[a,b)$ 上单调, 且 $\lim\limits_{x\to b-0}g(x)=0$,

则瑕积分 $\int_a^b f(x)g(x)\mathrm dx$ 收敛。

\noindent\textbf{定理 8.3.4（阿贝尔判别法）}\quad 设函数 $f(x),g(x)$ 在区间 $[a,b)$ 上满足下述条件:

(1) $\int_a^b f(x)\mathrm dx$ 收敛;

(2) $g(x)$ 在 $[a,b)$ 上单调有界,

则瑕积分 $\int_a^b f(x)g(x)\mathrm dx$ 收敛。

\noindent\textbf{例 8.3.5}\quad 讨论广义积分 $\int_0^{+\infty}\dfrac{\ln(1+x)}{x^\alpha}\mathrm dx(\alpha\in\mathbb R)$ 的敛散性。

\noindent\textbf{解}\quad 这个广义积分不仅是一个无穷积分, 而且当 $\alpha>1$ 时, 还有一瑕点 $x=0$. 为了讨论该广义积分的敛散性, 我们必须分别讨论
\[
  \int_0^1\frac{\ln(1+x)}{x^\alpha}\mathrm dx,\quad
  \int_1^{+\infty}\frac{\ln(1+x)}{x^\alpha}\mathrm dx
\]
的敛散性。

当 $x\to0+0$ 时, 我们有
\[
  \frac{\ln(1+x)}{x^\alpha}\sim\frac1{x^{\alpha-1}}.
\]
因此, 当 $\alpha<2$ 时, $\int_0^1\dfrac{\ln(1+x)}{x^\alpha}\mathrm dx$ 收敛; 当 $\alpha\geq2$ 时, $\int_0^1\dfrac{\ln(1+x)}{x^\alpha}\mathrm dx$ 发散。

当 $x\to+\infty$ 时, 如果 $\alpha\leq1$, 则有
\[
  \lim_{x\to+\infty}
  \frac{\dfrac{\ln(1+x)}{x^\alpha}}{\dfrac1{x^\alpha}}
  =
  \lim_{x\to+\infty}\ln(1+x)=+\infty.
\]
由于 $\int_1^{+\infty}\dfrac{\mathrm dx}{x^\alpha}$ 发散, 因此有 $\int_1^{+\infty}\dfrac{\ln(1+x)}{x^\alpha}\mathrm dx$ 发散。

当 $\alpha>1$ 时, 取 $\alpha'$, 使得 $1<\alpha'<\alpha$, 则有
\[
  \lim_{x\to+\infty}
  \frac{\dfrac{\ln(1+x)}{x^\alpha}}{\dfrac1{x^{\alpha'}}}
  =
  \lim_{x\to+\infty}\frac{\ln(1+x)}{x^{\alpha-\alpha'}}=0.
\]
由于 $\int_1^{+\infty}\dfrac{\mathrm dx}{x^{\alpha'}}$ 收敛, 从而 $\int_1^{+\infty}\dfrac{\ln(1+x)}{x^\alpha}\mathrm dx$ 收敛。

因此, 当 $1<\alpha<2$ 时, $\int_0^{+\infty}\dfrac{\ln(1+x)}{x^\alpha}\mathrm dx$ 收敛; 当 $\alpha\leq1$ 或 $\alpha\geq2$ 时, $\int_0^{+\infty}\dfrac{\ln(1+x)}{x^\alpha}\mathrm dx$ 发散。

\section*{习题八}

1. 计算下列无穷积分:

(1) $\int_1^{+\infty}\dfrac{\ln x}{(1+x)^2}\mathrm dx$;

(2) $\int_0^{+\infty}\dfrac{x\mathrm dx}{(x^2+a^2)^{\frac32}}\ (a\ne0)$;

(3) $\int_1^{+\infty}\dfrac{\mathrm dx}{x\sqrt{x^2+x+1}}$;

(4) $\int_1^{+\infty}\dfrac{\arctan x}{x^2}\mathrm dx$;

(5) $\int_0^{+\infty}\dfrac{\mathrm dx}{(1+x^2)^n}$。

2. 讨论下列无穷积分的敛散性:

(1) $\int_0^{+\infty}\dfrac{\sin x}{1+e^{-x}}\mathrm dx$;

(2) $\int_0^{+\infty}\dfrac{x^2+100x+1000}{x^4-x^3+1}\mathrm dx$;

\end{document}
